%-------------------------
% Resume in LateX
% Author : Gabriel Gomez
% License : MIT
%------------------------

\documentclass[letterpaper,11pt]{article}

\usepackage{cmap}
\usepackage[utf8]{inputenc}
\usepackage[T1]{fontenc}
\usepackage{lmodern}

\usepackage{latexsym}
\usepackage[empty]{fullpage}
\usepackage{titlesec}
\usepackage{marvosym}
\usepackage[usenames,dvipsnames]{color}
\usepackage{verbatim}
\usepackage{enumitem}
\usepackage[hidelinks]{hyperref}
\usepackage{fancyhdr}
\usepackage[spanish]{babel}
\usepackage{tabularx}
\input{glyphtounicode}

\pagestyle{fancy}
\fancyhf{} % Clear all header and footer fields
\fancyfoot{}
\renewcommand{\headrulewidth}{0pt}
\renewcommand{\footrulewidth}{0pt}

% Adjust margins
\addtolength{\oddsidemargin}{-0.5in}
\addtolength{\evensidemargin}{-0.5in}
\addtolength{\textwidth}{1in}
\addtolength{\topmargin}{-.5in}
\addtolength{\textheight}{1.0in}

\urlstyle{same}

\raggedbottom
\raggedright
\setlength{\tabcolsep}{0in}

% Sections formatting
\titleformat{\section}{
  \vspace{-4pt}\scshape\raggedright\large
}{}{0em}{}[\color{black}\titlerule \vspace{-5pt}]

% Ensure that generate PDF is machine readable/ATS parsable
\pdfgentounicode=1

%-------------------------
% Custom commands
\newcommand{\resumeItem}[2]{
  \item\small{
    \textbf{#1}{: #2 \vspace{-2pt}}
  }
}

% Just in case someone needs a heading that does not need to be in a list
\newcommand{\resumeHeading}[4]{
    \begin{tabular*}{0.99\textwidth}[t]{l@{\extracolsep{\fill}}r}
      \textbf{#1} & #2 \\
      \textit{\small#3} & \textit{\small #4} \\
    \end{tabular*}\vspace{-5pt}
}

\newcommand{\resumeSubheading}[4]{
  \vspace{-1pt}\item
    \begin{tabular*}{0.97\textwidth}[t]{l@{\extracolsep{\fill}}r}
      \textbf{#1} & #2 \\
      \textit{\small#3} & \textit{\small #4} \\
    \end{tabular*}\vspace{-5pt}
}

\newcommand{\resumeSubSubheading}[2]{
    \begin{tabular*}{0.97\textwidth}{l@{\extracolsep{\fill}}r}
      \textit{\small#1} & \textit{\small #2} \\
    \end{tabular*}\vspace{-5pt}
}

\newcommand{\resumeSubItem}[2]{\resumeItem{#1}{#2}\vspace{-4pt}}

\renewcommand{\labelitemii}{$\circ$}

\newcommand{\resumeSubHeadingListStart}{\begin{itemize}[leftmargin=*]}
\newcommand{\resumeSubHeadingListEnd}{\end{itemize}}
\newcommand{\resumeItemListStart}{\begin{itemize}}
\newcommand{\resumeItemListEnd}{\end{itemize}\vspace{-5pt}}
\setlength{\footskip}{5pt}
%-------------------------------------------
%%%%%%  CV STARTS HERE  %%%%%%%%%%%%%%%%%%%%%%%%%%%%


\begin{document}

%----------HEADING-----------------
\begin{tabular*}{\textwidth}{l@{\extracolsep{\fill}}r}
  \textbf{\href{https://gabogomez.dev/}{\Large Gabriel Gomez}} & Correo electrónico: \href{mailto:ggomez.dev@outlook.com}{ggomez.dev@outlook.com}\\
  \href{https://gabogomez.dev/}{gabogomez.dev} & Movil: \href{tel:+573168325201}{+57-316-832-5201} \\
\end{tabular*}


%-----------EXPERIENCE-----------------
\section{Experiencia}
\resumeSubHeadingListStart

\resumeSubheading
{Autónomo}{Barranquilla, CO}
{Desarrollador de Software Freelance}{Jun 2025 -- Presente}
\resumeItemListStart
\resumeItem{Migración de Web Service a Spring Boot}
{Migré un web service de Java EE a Spring Boot, rediseñando la capa de servicios y exponiendo APIs REST con mejor rendimiento y mantenibilidad. Implementé Kubernetes para la orquestación de contenedores y construí un entorno de desarrollo automatizado con Makefile que incluye un stack de Docker Compose con base de datos de pruebas, build de imágenes e inicialización de configuración.}
\resumeItemListEnd

\resumeSubheading
{Extreme Technologies S.A.}{Barranquilla, CO}
{Desarrollador de Software}{Abr 2021 -- Jun 2024}
\resumeItemListStart
\resumeItem{Acuacar Daños}
{Mantuve y desarrollé funcionalidades en un web service Jakarta EE que expone APIs REST consumidas por una app Android para gestión de cuadrillas en la nube. Rediseñé la interfaz de la app móvil, migré el soporte a Android 11 e implementé un flujo de trabajo que abarcó tanto el backend como el cliente móvil, coordinando los cambios de contrato de API entre ambas capas.}
\resumeItem{Promigas}
{Adapté un sistema white-label de gestión de cuadrillas para un cliente del sector gas, desarrollando el web service Jakarta EE y la app Android cliente. Implementé rastreo GPS en tiempo real integrando PostGIS para el procesamiento de coordenadas geográficas, y optimicé el backend para gestionar la carga incrementada por sesiones concurrentes en tiempo real, mejorando la estabilidad bajo alta demanda.}
\resumeItem{Podas y Lavados}
{Adapté un sistema white-label de mantenimiento de redes eléctricas para los clientes Aire y Afinia sobre Jakarta EE. Desarrollé un feature GIS en el backend que descarga y sirve capas con posiciones y estados de árboles y postes eléctricos, y construí la app Android multirol con soporte para dos flujos de trabajo paralelos que consumen estas APIs.}
\resumeItem{SETP Santa Marta}
{Implementé una capa middleware en un validador de pasajes Android para el sistema de transporte público de Santa Marta. Desarrollé la lógica de validación de pagos para tarjetas físicas y códigos QR de billeteras virtuales, y diseñé una cola offline-first para almacenar transacciones durante pérdidas de conexión y sincronizarlas al restaurar la conectividad.}
\resumeItemListEnd

\resumeSubHeadingListEnd

%-----------EDUCATION-----------------
\section{Educación}
\resumeSubHeadingListStart
\resumeSubheading
{Universidad Autónoma del Caribe}{Barranquilla, CO}
{Pregrado en Ingeniería de Sistemas}{Feb 2018 -- May 2023}
\resumeSubHeadingListEnd


%-----------PROJECTS-----------------
\section{Proyectos}
\resumeSubHeadingListStart
\resumeSubItem{MatricuApp}
{Webapp de matrícula universitaria con backend en Spring Boot, Maven, JPA e Hibernate, autenticación JWT con Spring Security y base de datos PostgreSQL. Expone APIs REST para que los estudiantes inicien sesión, matriculen y retiren materias con validación de límite de créditos. Frontend construido con React.}
\resumeSubItem{Portfolio Personal}
{Portfolio web construido con Astro y Tailwind CSS, desplegado en Cloudflare. Implementa Edge Functions para redirección automática basada en el idioma del navegador, ofreciendo una experiencia multilingüe de alto rendimiento.}
\resumeSubItem{Resume Builder en LaTeX}
{Generador de currículos de código abierto con LaTeX y GitHub Actions para compilación y despliegue automatizados de PDFs, siguiendo prácticas de CI/CD.}
\resumeSubHeadingListEnd

%-----------Cursos y Certificaciones-----------------
\section{Cursos y Certificaciones}
\resumeSubHeadingListStart
\resumeSubItem{Udemy - Curso de React: de cero a experto}{Mar 2026}
\resumeSubItem{Platzi - Herramientas de AI para Developers}{Abr 2025}
\resumeSubItem{Udemy - Curso Testing para Android con JUnit, Mockito, Espresso y TDD}{Feb 2024}
\resumeSubItem{Universidad Autónoma del Caribe - Nuevas Tendencias de Ingeniería}{Feb 2023}
\resumeSubItem{Platzi - Java SE Orientado a Objetos}{Feb 2021}
\resumeSubItem{Platzi - Introducción a Java SE}{Oct 2020}
\resumeSubItem{Platzi - Curso Profesional de Git y GitHub}{Sep 2019}
\resumeSubHeadingListEnd

%--------SKILLS------------
\section{Skills}
\resumeSubHeadingListStart
\item{
            \textbf{Lenguajes}{: Java, SQL, TypeScript, JavaScript, Kotlin, Python, HTML, CSS, LaTeX}
      }
\item{
            \textbf{Tecnologías}{: Spring Boot, Jakarta EE (Java EE), REST APIs, PostgreSQL, MySQL, JUnit, Mockito, Test-Driven Development, Docker, Kubernetes, Docker Compose, GitHub Actions, CI/CD, Clean Architecture, SOLID, Design Patterns, Microservices, Maven, Gradle, Agile, Scrum, Jira, Git, Firebase}
      }
\resumeSubHeadingListEnd


%-------------------------------------------
\end{document}
